\documentclass{ximera}
% xmPreamble.tex en xmPrintstyle.sty worden automatisch geladen door ximera.cls (v2.7.8);
% NIET nog eens \input{./preamble.tex} of \addPrintStyle{.} doen (geeft 'already defined'-fouten).

\begin{document}
\author{Alexander Holvoet}
\xmtitle{Een oefening op convexe veelhoeken}{}

\begin{exercise}
    Construeer, door gebruik te maken van bissectrices, de bovenaanzichten van de zandhopen
    met de gegeven grondvlakken. Je kunt hiervoor GeoGebra gebruiken. Verifieer daarna met
    een experiment.
    \begin{hint}
        De kamlijn bestaat uit de punten die even ver liggen van de twee dichtstbijzijnde
        randen. Dat zijn net stukken bissectrice tussen opeenvolgende zijden.
    \end{hint}
    \begin{oplossing}
        Teken bij elke veelhoek de bissectrices van de hoeken; de kamlijnen vind je waar de
        bissectrices van opeenvolgende zijden samenkomen. De aardrijkskundige hoogtelijnen
        helpen om de constructie te controleren, maar zijn niet nodig om de oefening op te
        lossen.
    \end{oplossing}
\end{exercise}

% TODO (afbeeldingen ontbreken): in de bron stonden hier Lconvexeveelhoeken.png (de op te
% lossen grondvlakken) en Lconvexemethoogtel.png (de oplossingen met hoogtelijnen). Niet
% aanwezig in img/.

\end{document}
